The issue of ensuring safety for aerial robots when navigating near obstacles remains a significant concern. There are three distinct aspects to addressing these safety concerns: the controller, which must provide assurances of safety while following a prescribed trajectory; the trajectory planner, which must generate a trajectory that avoids all obstacles; and the perception module, which must accurately localize the robot and obstacles with minimal deviation from actual ground truth. This research study focuses specifically on the first aspect, which pertains to the safe control of aerial robots when maneuvering near static obstacles. The study considers two types of aerial robots, namely an aerial manipulator and a quadrotor unmanned aerial vehicle.


The AM considered here is a 2-DOF (degrees-of-freedom) manipulator rigidly attached to a UAV. Our proposed controller structure follows the conventional inner loop PID control for attitude dynamics and an outer loop controller for tracking a reference trajectory. The outer loop control is based on the Model Predictive Control (MPC) with constraints derived using the Barrier Lyapunov Function (BLF) for the safe operation of the AM. BLF-based constraints are proposed for two objectives, viz. 1) To avoid the AM from colliding with static obstacles like a rectangular wall, and 2) To maintain the end effector of the manipulator within the desired workspace.
The proposed BLF ensures that the above-mentioned objectives are satisfied even in the presence of unknown bounded disturbances. The capabilities of the proposed controller are demonstrated through high-fidelity non-linear simulations with parameters derived from a real laboratory scale AM. We compare the performance of our controller with other state-of-the-art MPC controllers for AM.

 To address the problem of flying a UAV inside a tunnel, we propose the implementation of a Model Predictive Control (MPC) framework with constraints based on Control Barrier Function (CBF). The thesis approaches the issue in two distinct ways; first, by maintaining a safe distance from the tunnel walls to avoid the effects of both the walls and ceiling, and second, by minimizing the distance from the walls to effectively manage the nonlinear forces associated with close proximity tasks. Finally, the paper demonstrates the effectiveness of its approach through testing on simulation for various close proximity trajectories with the realistic model of aerodynamic disturbances due to the proximity of the ceiling and boundary walls.

 